\subsubsection{A complex rank-one singular set: weights \texorpdfstring{$(3,7,7)$}{(3,7,7)}}
\label{sec:weighted-Clifford-quartic}

The balanced criterion \cref{prop:balanced-quartic-multiplier} gives
a harmonic example with a complex rank-one singular set.
Let $E_i=L_{e_i}$, $1\leq i\leq7$, be left multiplication by an
orthonormal imaginary basis of the octonions.  Alternativity gives
$E_i^T=-E_i$ and $E_iE_j+E_jE_i=-2\delta_{ij}\Id_8$.
On $\R^{16}=\R^8\oplus\R^8$ put
\begin{equation}\label{eq:weighted-Clifford-matrices}
 J=\begin{pmatrix}0&\Id_8\\-\Id_8&0\end{pmatrix},\qquad
 B_i=\begin{pmatrix}E_i&0\\0&-E_i\end{pmatrix}.
\end{equation}
These eight skew-symmetric matrices satisfy the negative Clifford
relations.  For $x,y\in\R^{16}$ define
\begin{align}
 R&=|x|^2+|y|^2,\nonumber\\
 W&=|x|^2|y|^2-\langle x,y\rangle^2-\langle Jx,y\rangle^2,
 \label{eq:weighted-Clifford-definition}\\
 F&=2\sum_{i=1}^7\langle B_ix,y\rangle^2-W.\nonumber
\end{align}

\begin{proposition}[The complex rank-one example]
\label{thm:weighted-Clifford-quartic}
The polynomial $F$ in \eqref{eq:weighted-Clifford-definition} is
real irreducible and harmonic. Its cone $C=\{F=0\}\subset\R^{32}$ has
geometric singular set
\begin{equation}\label{eq:weighted-Clifford-singular-set}
 \Sing C=\Sigma=\{W=0\}
 \cong\{(z,w)\in\mathbb C^8\times\mathbb C^8:
             \operatorname{rank}_{\mathbb C}[z\ w]\leq1\}.
\end{equation}
This set has real dimension $18$ and is smooth away from the origin.
At each nonzero singular point the tangent cone is orthogonally
congruent to $\R^{18}\times\{|\xi|^2=|\eta|^2\}$, with
$\xi,\eta\in\R^7$. The potential $G=WF$ subcalibrates the full
boundary current, which is strictly area minimizing. Its smallest
Laplacian algebra is $\R[R,W,F]$, and \eqref{eq:bihomogeneous-Hardy} gives
\begin{equation}\label{cor:weighted-Clifford-Hardy}
 Q_C(\varphi)\geq\frac{385}{4}\int_C r^{-2}\varphi^2,
 \qquad \varphi\in C_c^\infty(C_{\reg}\setminus\{0\}).
\end{equation}
\end{proposition}

\begin{proof}
The Clifford system associated with $\Id_{16},J,B_1,\ldots,B_7$
has quadratic coordinates $p_0=|x|^2-|y|^2$ and
$p_A=2\langle Ax,y\rangle$. Apply \cref{lem:Clifford-block-quartics} with
\begin{align}
 E_0&=p_0^2+p_{\Id}^2+p_J^2=R^2-4W,\nonumber\\
 E_+&=\sum_{i=1}^7p_{B_i}^2=2(W+F),\label{eq:weighted-Clifford-blocks}\\
 E_-&=R^2-E_0-E_+=2(W-F).\nonumber
\end{align}
Here $\ell=16$, $m=9$, and the weights are $(3,7,7)$, so
\cref{prop:balanced-quartic-multiplier} gives harmonicity, the smallest
algebra, and the subcalibration. Irreducibility follows from
\cref{prop:minimizing-irreducibility}; the Hardy coefficient uses
$(N,h)=(32,14)$.

Identify $(\R^{16},J)$ with $\mathbb C^8$.  The expression for
$W$ is the complex Gram determinant of the two columns $x,y$.
Thus \cref{lem:quartic-table-geometry} and equality in Cauchy--Schwarz
give \eqref{eq:weighted-Clifford-singular-set}.
The rank-one matrices form a smooth complex manifold of dimension
$8+2-1=9$, so $\dim_{\R}\Sigma=18$ and
$\Sigma\setminus\{0\}$ is smooth.

For $x\ne0$, the vectors $x,Jx,B_1x,\ldots,B_7x$ are orthogonal
with length $|x|$. Decompose $y=y_0+y_++y_-$ along
$K_x=\operatorname{span}\{x,Jx\}$,
$U_x^+=\operatorname{span}\{B_1x,\ldots,B_7x\}$, and
$U_x^-=(K_x\oplus U_x^+)^\perp$. The last two spaces have dimension seven, and
\begin{equation}\label{eq:weighted-Clifford-transverse-form}
 W=|x|^2(|y_+|^2+|y_-|^2),\qquad
 F=|x|^2(|y_+|^2-|y_-|^2).
\end{equation}
Local frames give transverse signature $(7,7)$; interchanging $x,y$
covers $\Sigma\setminus\{0\}$. The Hessian formulas in
\cref{lem:quartic-table-geometry}, with $T=W$ and $F_{\rm gen}=F/2$,
give eigenvalues $\pm2R(p)$, each sevenfold, and kernel dimension
eighteen; \cref{sec:tangent-instability-transfer} gives the asserted
orthogonal tangent-cone model.
\end{proof}

The singular dimension $18$ distinguishes this cone from the
splitting and Hopf pairing cones in $\R^{32}$, whose singular sets
have dimension $16$. For a pure full Clifford pullback of a
nondegenerate Lawson cone, \eqref{eq:Clifford-pullback-identities}
makes the nonvertex singular set the zero fiber of the Clifford map.
Congruence would therefore require fourteen quadratic coordinates.
Such a system on $\R^{32}$ would require $\delta(13)=128\mid16$,
contrary to the module restriction \cite[Section~3.5]{FKM81}.
In particular, the example is distinct from Riedler's family.
